\begin{abcliste}
\abc Across the part of the fin between the magnets the motional emf is induced:
\begin{align}
 U &= \UiF = \Bm\times\hf\times\v = \Ui \approx \result{\UiP{0}{2}}
\end{align}

\abc By Lenz's rule, the currents induced in the fin (eddy currents) oppose the change that causes them: the motion of the fin through the field. The magnets $\result{\text{brake the fin}}$. (Copper is not attracted by magnets, but currents are induced in it.)

\abc For the same reason they $\result{\text{brake it}}$ while it leaves the field: the induced currents always oppose the motion, whether the fin moves into the field or out of it.

\abc No. At rest, the flux does not change: no current is induced, and there is $\result{\text{no force}}$. A magnetic brake can slow a car down, but it cannot hold it.
\end{abcliste}
